\cotefile{56}
\scanpage{28}
% archive sheet number: 27
% The table reads across the sheet when the scan is turned clockwise.
\diagram{\small\setlength{\arraycolsep}{2.4pt}
$\begin{array}{cccccccc}
\left[\begin{smallmatrix}1,1\\1,-1\end{smallmatrix}\right] &
 Q^{0}&Q^{1}&Q^{2}&Q^{3}&Q^{4}&Q^{5}&Q^{6}\\
&[\lambda_{0},\xi^{0}]&\xi^{0}=1&\xi^{0}=1&\xi^{0}=1&\xi^{0}=1&\xi^{0}=1&\xi^{0}=1\\
&&\frac{1}{2}\xi^{\deleted{\unsure{1|2}}}&[\lambda_{1},\xi]&\xi&\xi&\xi&\xi\\
&&&\frac{1}{2}\xi^{2}&\frac{1}{2}\xi^{2}&[\lambda_{2},\xi^{2}]&\xi^{2}&\xi^{2}\\
&&&&\frac{1}{2}\xi^{3}&\frac{1}{2}\xi^{3}&\frac{1}{2}\xi^{3}&[\lambda_{3},\xi^{3}]\\
&&&&&\frac{1}{2}\xi^{4}&\frac{1}{2}\xi^{4}&\frac{1}{2}\xi^{4}\\
&&&&&&\frac{1}{2}\xi^{5}&\frac{1}{2}\xi^{5}\\
&&&&&&&\frac{\unsure{1|+}}{2}\xi^{6}
\end{array}$}

% margin note written sideways
\margin{a) \unsure{Chgt pour|Changement pour|Cohomologie} de $Q^{n}$ nulle en dim impaire}
\scanpage{29}
% archive sheet number: 28
\mspage{\handwritten{12}}
% L1–3
Alors $T(k)$ est dans le centre de $G(k)$ en vertu de BIBLE 6. th.2 , ce qui implique que $T$ est dans le centre de $G$, \deleted{en vertu de Exp VI} \inserted{\handwritten{car $\underline{\mathrm{Centr}}_{G}(T)$ est un sous-schéma fermé de $G$ (Exp \uline{VIII} 6.6.), qui contient les points de $G(k)$, donc est identique à $G$ par le Nullstellensatz, ($G$ étant réduit).}}
\margin{\handwritten{Nullstellensatz}}

% L4–5
Enfin, pour référence ultérieure, signalons la proposition triviale suivante, (dont nous avons fait déjà usage implicitement):

% L6–7
\uline{Proposition} \handwritten{1}.17. \quad Soient $G\handwritten{\supset}H\handwritten{\supset}T$ de\handwritten{s} préschémas en groupes de type fini sur $S$. \deleted{Alors:} \handwritten{Conditions équivalentes:}

% L8–9
\handwritten{(i)} \deleted{\unsure{Si}} $T$ est un tore maximal de $G$, \deleted{$T$ est un tore maximal de $H$, \handwritten{et pour tout $s\in S$, \unsure{$H_s$ et $G_s$ ont même rang réductif|$G_s$ et $H_s$ ont même rang réductif}}}

% L10–12
\handwritten{(ii)} \deleted{\unsure{Si}} $T$ est un tore maximal de $H$, et \deleted{\unsure{si}} pour tout $s\handwritten{\in}S$, on a égalité des rangs réductifs $\handwritten{\rho}_{r}(G_{s})=\handwritten{\rho}_{r}(H_{s})$\handwritten{,} \deleted{alors $T$ est un tore maximal de $G$}.

% L13
\handwritten{2}. \quad \uline{Le groupe de Weyl}.

% L14–19
Soit d'abord $k$ un corps algébriquement clos, et soit $G$ un groupe algébrique sur $k$, \inserted{lisse et affine sur $k$.} Si $T$ est un tore maximal, $C$ son centralisateur et $N$ son normalisateur, alors en vertu de \inserted{\handwritten{Exp \uline{XI}}}5.9. ce sont des sous-groupes fermés \deleted{\unsure{essentiellement}} \inserted{lisses de $G$,} et $\handwritten{C}$ est un sous-groupe ouvert de $N$, de sorte que $\inserted{W=}N/C$ est un groupe fini \deleted{\unsure{étale, qu'on appelle le groupe de Weyl}} \inserted{étale sur $k$ , donc déterminé par le groupe $W(k)$} de \deleted{\unsure{$G$ relatif au tore maximal $T$.}} ses points à valeurs dans $k$, comme $W=W(k)_{k}$ (groupe constant défini par le groupe fini ordinaire $W(k)$).
% L20–25
Le groupe fini $W(k)$ s'appelera le \uline{groupe de Weyl géométrique}, \handwritten{(}ou simplement groupe de Weyl si une confusion n'est pas à craindre\handwritten{)}, \uline{de} $G$ \uline{relativement} \uline{à} $T$. En vertu du théorème de \inserted{conjug\deleted{u}aison \handwritten{de}} Borel, les groupes de Weyl relatifs aux divers tores maximaux sont isomorphes entre eux, c'est pourquoi on parle parfois ``du'' groupe de Weyl de $G$, sans préciser de tore maximal. \deleted{\unsure{En}}

% L26–28
\handwritten{C}omme la formation de $C,N$, et $N/C$ commute à toute extension de la base, on voit que si $k'$ est une extension algébriquement close de $k$, le groupe de Weyl géométrique de $G_{k'}$ relativement à $T_{k'}$\handwritten{\unsure{/}} est canoniquement isomorphe
\scanpage{30}
% archive sheet number: 29
% L1
\[
\deleted{\boxed{\mathord{\text{\textcircled{\ensuremath{F}}}}\otimes G}}
\]
% L2
\[
L_{\cdot}(F)\otimes L_{\cdot}(G)=(F\mathbin{\underline{\underline{\otimes}}}G)
\]

% L3–11
\diagram{\scriptsize
\begin{tikzcd}[ampersand replacement=\&,column sep=small,row sep=small,
 squiggly/.style={decorate,decoration={snake,amplitude=.55pt,segment length=4pt,post length=1.4pt}}]
{} \& H^{n+1}(\unsure{U|W}^{n+1})(1)
 \arrow[r,"\mathrm{inj}" ,"\partial"'] \arrow[rr,bend left=28,"{\substack{\deleted{\unsure{\mathrm{bij}}}\,\deleted{\unsure{8}}\,\mathrm{bij}\ \underline{\mathrm{exc}\ 2},\\ \ell^{(2\nu+1)}\rightsquigarrow2\ell^{\overset{\unsure{\curvearrowleft}}{(2\nu)}}}}"]
 \& H^{n}(Q^{n})\arrow[r,"\mathrm{surj}"] \& H^{n}(U^{n})
 \& {}\arrow[ddddddd,"{n=2\nu}"{pos=.04}]\\
{} \& \ell_{\deleted{\unsure{2\nu+1}}}^{2\nu+1} \arrow[l,shorten <=15pt]
 \& {[\lambda_{\nu},\xi^{\nu}]}\arrow[r,squiggly] \& {[2\ell_{\deleted{\unsure{2\nu}}}^{2\nu},0]} \&\\
\& {\ell^{2\nu+1}\rightsquigarrow\lambda_{\nu}}
 \& {\varphi_{\nu}=\frac12(\lambda_{\nu}+\xi^{\nu}\unsure{]|)}}\arrow[r,squiggly]
 \& \ell_{\deleted{\unsure{2\nu}}}^{2\nu} \&\\
{}\arrow[rrrr,dash] \& {} \& \deleted{H^{0}(Q^{0})}\arrow[r]
 \& \deleted{H^{0}(U^{\unsure{\nu|0}})} \& {}\\
\& \& \lambda_{\nu} \& \ell_{2\nu}\arrow[l,squiggly] \&\\
\& H^{n+1}_{c}(U^{n+1})(+1)
 \& H^{\unsure{n}}(Q^{n})\arrow[l,"\mathrm{surj}"',"\partial"]
 \& H^{n}_{c}(U^{n})\arrow[l,"\mathrm{inj}"'] \&\\
\& {[2\ell_{2\nu+1},0]}\& {[\lambda_{\nu},\xi^{\nu}]}\arrow[l,squiggly] \& \& {}\\
\& \ell_{2\nu+1} \& \varphi_{\nu}\arrow[l,squiggly] \& {}\arrow[ll,bend left=22] \& {}\\
\& 2\ell_{2\nu+1} \& {} \& \ell_{2\nu}\arrow[ll,squiggly] \& {}
\end{tikzcd}}

% L12
\[
Q^{n+1}-Q^{n}=U^{n+1}
\]

% L13–17
\uline{N.B.} \quad \uline{On définit} $\deleted{\unsure{\lambda}}\varphi_{\nu}$ \uline{directement} \uline{par les sous-espaces isotropes maximaux.} -- les \unsure{,} deux familles de tels sous-espaces \unsure{correspondent} \unsure{:|à} des classes $\alpha_{\nu},\beta_{\nu}$. On a $\alpha_{\nu}+\beta_{\nu}=\xi^{\nu}$, $\alpha_{\nu}-\beta_{\nu}=\lambda_{\nu}$, $\varphi_{\nu}=\alpha_{\nu}$ \unsure{---|\ldots}
\scanpage{31}
% archive sheet number: 30
\mspage{\handwritten{26}}
% L1
Démonstration.

% L2–8
a) et b). La démonstration de a) et de la première assertion dans b) est essentiellement identique à celle de 1.7. a) et b), et nous nous dispensons de reproduire ici le raisonnement. Signalons seulement qu'il faut faire appel dans a) à Exp IX 5.6. (utilisant le fait que $G$ est à fibres connexes). Prouvons la deuxième assertion de b), i.e. que si $Z$ est un centre réductif de $G$, alors $G'=G/Z$ admet pour centre réductif le sous-groupe unité.
% L8–15
Notons tout de suite qu'en vertu de Exp VIII 5.1. le groupe quotient $G/Z$ estb bien représentable, il est affine sur $S$, de présentation finie sur $S$ (Exp VIII 5.8.) et même lisse sur $S$ (car il suffit de le voir pour les fibres, $G'$ étant plat de présentation finie sur $S$, or pour les fibres \deleted{cela est} on a signalé dans Exp VI qu'un quotient d'un groupe algébrique lisse sur un corps $k$ est lisse); enfin $G$ étant à fibres connexes, il en est de même de $G'$. Ainsi $G'$ satisfait aux mêmes hypothèses de départ que $G$.
% L15–21
Pour voir que le sous-groupe unité de $G'$ est un centre réductif, on peut se borner si on veut au cas où $\deleted{\unsure{s}}\handwritten{S}$ est le spectre d'un corps (4.\deleted{\unsure{2|3}}\handwritten{3}.). Soit $Z'$ un sous-groupe central de type multiplicatif de $G'$, tout revient à prouver que $\handwritten{Z}'$ est réduit au groupe unité. Soit $Z_{1}$ \deleted{\unsure{l'im}} l'image inverse de $Z'$ dans $G$, et considérons les opérations de $Z_{1}$ sur $G$ induites par automorphismes intérieurs de $G$. Comme $Z$ est central, $Z$ opère \deleted{\unsure{tri}} trivialement, de sorte que $Z_{1}$ opère par l'intermédiaire d'opérations de $Z'$ sur $G$.
% L22–26
D'ailleurs $Z$ étant dans le centre de $G$, $Z\deleted{\unsure{'}}_{1}$ donc $Z'$ opère trivialement sur $Z$, de plus $Z'$ étant dans le centre de $G'$, les opérations de $Z'$ dans le quotient $G/Z=G'$ \handwritten{\unsure{/|(}}sont également triviales. Comme $Z$ est central, il s'ensuit aussitôt que les opérations de $Z'$ sur $G$ proviennent d'un
\scanpage{32}
% archive sheet number: 31
% L1
\[
Q^{2\nu}:\qquad \xi:\quad
H^{2i}(\deleted{\unsure{X_{2}|X_{\mathbb Z}}}Q^{2\nu})\longrightarrow
H^{2i+2}(\deleted{\unsure{X_{2}|X_{7}}}Q^{2\nu})(\unsure{i|1})
\]
% L2–5
\[
\begin{aligned}
&\deleted{\text{\uline{injectif} si }2i<2\nu\unsure{<}\ \text{i.e. }i\leq\nu-1}\\
&\text{\uline{bijectif} si }\ i\ne\nu-1,\nu
\qquad \deleted{\unsure{(i\geq0)}}\ (0\leq i\leq2\nu-1)\\
i=\nu-1:\quad &\text{\uline{injectif} }\deleted{\text{si }i=\nu\unsure{{-}1}},\
 \text{conoyau libre de base }\varphi_{\nu}\\
i=\nu\quad &\deleted{\text{\unsure{noyau}}}\text{\inserted{\uline{surjectif}}},\
 \text{noyau libre de base }\lambda_{\nu}
\end{aligned}
\]
\noindent\rule{.40\linewidth}{.4pt}

% L6–8
\[
\begin{aligned}
Q^{2\nu+1}\qquad
&\unsure{\xi}\quad H^{2i}(Q^{2\nu+1})\unsure{-|\longrightarrow}H^{\unsure{2i+2|2i+1}}(Q^{2\nu\inserted{+1}})(1)\\
&\text{\uline{bijectif} si }\ i\ne\nu\qquad (0\leq i\leq2\nu)\\
&\text{injectif de conoyau }\mathbb Z/2\mathbb Z\quad\text{si }i=\nu
\end{aligned}
\]
\noindent\rule{.83\linewidth}{.4pt}
\scanpage{33}
% archive sheet number: 32
\mspage{\handwritten{31}}
% L1–3
En effet, on peut supposer $k$ algébriquement clos, de sorte que $T$ est contenu dans un tore maximal \deleted{\unsure{de $G$}} $T'$. Soit $C'$ le centralisateur de ce dernier, alors on a $C'\handwritten{\subset}C$, donc $\dim C'\handwritten{\leq}\dim C$, cqfd.

% L4–8
b) \deleted{\unsure{Le raisonnement de a) nous montre que}} Si $\handwritten{\rho}_{r}$ est locqlement constant, alors pour tout tore $T$ dans $G$, \deleted{\unsure{et}} et tout $s\handwritten{\in}S$, si $T_{s}$ est \inserted{tore} maximal dans $G_{s}$, alors il existe un voisinage ouvert $U$ de $s$ tel que $\deleted{\unsure{G}}T\handwritten{\vert}U$ soit tore maximal dans $G\ U$. Utilisant maintenant le raisonnement de a), on voit que \deleted{\unsure{si}} (i)$\handwritten{\Rightarrow}$(ii bis). D'autre part (ii bis)$\handwritten{\Rightarrow}$(i),
% L9–15
car \deleted{on peut supposer $S$} lorsque $G$ admet un tore maximal $T$, alors il est évident que $\handwritten{\rho}_{r}(s)=\dim T_{s}$ est une fonction localement constante de \deleted{\unsure{$s$}} $s$, or nous avons signalé au début de la démonstration de a) que la question de la continuité de $\handwritten{\rho}_{r}$ était locale pour la topologie fidèlement plate. D\deleted{\unsure{v}}onc (i)$\handwritten{\Leftrightarrow}$(ii bis), et évidemment (ii)$\handwritten{\Rightarrow}$(ii bis), reste à prouver l'implication inverse \inserted{(i)$\handwritten{\Longrightarrow}$(ii)}.
% L15–19
Pour ceci, introduisons le foncteur $F$ de \inserted{\handwritten{Exp \uline{XI}}}4.1. , qui est donc un préschéma \deleted{\unsure{localement de}} lisse \inserted{\handwritten{et séparé}} sur $S$, et considérons le sous-foncteur $\deleted{\unsure{F}}\handwritten{\mathscr{T}}$ de $F$, dont la valeur \deleted{sur} pour un $\handwritten{S'}$ sur $S$ est l'ensemble des tores maximaux dans $G_{\handwritten{S'}}$. Utilisant la remarque précédente \deleted{\unsure{que}} \inserted{que\handwritten{,} moyennant (i)\handwritten{,}} un tore dans ${G\deleted{\unsure{s|a}}}_{\handwritten{S'}}$ qui est maximal dans la fibre d'un
% L20–24
point $\handwritten{s\in S'}$ est maximal au dessus d\deleted{\unsure{e}}'un voisinage ouvert de $\handwritten{s}$, on voit que $\handwritten{\mathscr{T}}$ \margin{\handwritten{$\mathscr{T}$}}est représentable par un sous-préschéma ouvert de $F$, et par suite est lisse \inserted{\handwritten{et séparé}} sur $S$. Comme le morphisme structural $\handwritten{\mathscr{T}\to S}$ est évidem\deleted{\unsure{-}}ment surjectif, il admet donc \deleted{localement} une section localement au sens de la topologie étale en vertu de \inserted{\handwritten{Exp \uline{XI}}}1.1\handwritten{0}. , ce qui \handwritten{p}rouve (i)$\handwritten{\Rightarrow}$(ii).

% L25–26
c) C'est une conséquence immédiate de \inserted{\handwritten{Exp \uline{XI}}}5.4. bis, compte tenu du théorème de conjugaison de Borel rappelé au début du numéro.

% L27–28
d) Moyennant la remarque du début de la démonstration dans a), et compte tenu de b), on peut supposer qu'il existe un tore maximal $T$
\scanpage{35}
% archive sheet number: 34
% L1–3
\uline{Quadriques} \quad Une base pour l'équivalence linéaire dans $Q^{2n}$ est \unsure{formé|formée} de
\[
1,H,H^{\unsure{2|7}},\ldots H^{n-1},\
\underbrace{S',S''}_{\text{codim }n},\
HS'=HS'',\ \ldots\ H^{n}S'=H^{n}S''
\]
\margin{On a $H^{n}=S'+S''$}

% L4–7
[(\,$H$ est une section hyperplane dans l'immersion projective évidente; $S'$ et $S''$ sont définis par des sous-espaces \deleted{\unsure{affines}} \inserted{linéaires de dim $n$} dans $Q^{2n}$, appartenant aux deux familles algébriques distinctes de tels sous-espaces\,]

% L8–9
Une base de $Q^{2n+1}$ est formée de
\[
1,H,\ldots H^{n},\frac{H^{n+1}}{2},\ldots,\frac{H^{2n+1}}{2}
\]
% L10–12
[où $H$ est une section hyperplane, et où les \unsure{termes|formes} $\frac{H^{n+1}}{2}\ldots\frac{H^{2n+1}}{2}$ correspondent aux sous-variétés linéaires contenues dans $Q^{2n+1}$\,]

\begin{center}\rule{.23\linewidth}{.4pt}\\[-.65em]\rule{.17\linewidth}{.4pt}\end{center}
\scanpage{36}
% archive sheet number: 35
% L1–3
\begin{center}\rule{.88\linewidth}{.35pt}\\[-.85em]\rule{.88\linewidth}{.65pt}\end{center}
\textsc{Algèbre}. --- \emph{Définition des classes de Stiefel-Whitney d'un module quadratique sur un corps de caractéristique différente de $2$.} Note de M. \textbf{\textsc{Antoine Delzant}}, transmise par M. Jean Leray.

% L4–5
Soit $\mathrm A$ un corps de caractéristique différente de $2$. On ne considère que des modules quadratiques $(\mathrm E,\mathrm Q)$ $(^{1})$, de rang fini sur $\mathrm A$, et non dégénérés.

% L6–16
On désigne par $\mathrm Q^{+}$ l'ensemble des classes à isomorphisme près de modules quadratiques. L'ensemble $\mathrm Q^{+}$ est muni de deux opérations : l'une, la somme directe externe, définit sur $\mathrm Q^{+}$ une structure de monoïde abélien pour laquelle tous les éléments sont réguliers (théorème de Witt); l'autre, le produit tensoriel, est commutative, distributive par rapport à la précédente, et admet un élément unité [la classe du module $(\mathrm A,\mathrm I)$, où $\mathrm I(x)=x^{2}$]. L'ensemble $\mathrm Q^{+}$ se plonge donc injectivement dans un anneau $\mathrm Q(\mathrm A)$, qu'il engendre. L'anneau $\mathrm Q(\mathrm A)$, commutatif avec élément unité, est appelé l'\emph{anneau des classes de formes quadratiques}. A toute classe $q$ de module quadratique $(\mathrm E,\mathrm Q)$ on fait correspondre $\varepsilon(q)=rg\,(\mathrm E)$; on définit ainsi un homomorphisme d'anneau de $\mathrm Q(\mathrm A)$ dans $\mathbf Z$ appelé l'augmentation.

% L17–21
Soit $\mathrm A_{s}$ la clôture séparable de $\mathrm A$ et $\mathrm G$ le groupe de Galois de $\mathrm A_{s}$ sur $\mathrm A$; on note $\mathrm H^{q}(\mathrm A,\mathbf Z/2\,\mathbf Z)$ le $q^{\text{--ième}}$ groupe de cohomologie de $\mathrm G$ à valeurs dans $\mathbf Z/2\,\mathbf Z$. On note additivement le groupe $\mathrm A^{\star}$, groupe multiplicatif de $\mathrm A$, et l'on pose $\mathrm D(\mathrm A)=\mathrm A^{\star}/2\,\mathrm A^{\star}$. Le groupe $\mathrm D(\mathrm A)$ s'identifie à $\mathrm H^{1}(\mathrm A,\mathbf Z/2\,\mathbf Z)$.

% L22–27
Soit $a$ un élément de $\mathrm A^{\star}$; on note $e_{a}$ sa classe dans $\mathrm D(\mathrm A)$, et $\varphi_{a}$ la classe du module quadratique $[\mathrm A,\mathrm Q;\ \mathrm Q(x)=ax^{2}]$. L'élément $\varphi_{a}$ ne dépend que de $a$ modulo $2\,\mathrm A^{\star}$; on peut donc définir un homomorphisme de $\mathrm D(\mathrm A)$ dans le groupe multiplicatif $\mathrm Q(\mathrm A)^{\star}$ des éléments inversibles de $\mathrm Q(\mathrm A)$. D'où un homomorphisme d'anneaux :
\[
\varphi:\quad \mathbf Z(\mathrm D)\to\mathrm Q(\mathrm A),
\]
% L28–29
où $\mathbf Z(\mathrm D)$ est l'algèbre de $\mathrm D(\mathrm A)$ sur $\mathbf Z$, munie de la filtration définie par les puissances de l'idéal d'augmentation.

% L30–32
\textsc{Proposition}. --- \emph{L'application $\varphi$ est surjective; elle est compatible avec les augmentations; le noyau de $\varphi$ est l'idéal engendré par les éléments $e_{a}+e_{b}-e_{a'}-e_{b'}$, tels que $\varphi_{a}+\varphi_{b}=\varphi_{a'}+\varphi_{b'}$.}

% L33–35
Le cup-produit définit une application de $\mathrm D(\mathrm A)\times\mathrm D(\mathrm A)$ dans $\mathrm H^{2}(\mathrm A,\mathbf Z/2\,\mathbf Z)$ notée $(e_{a},e_{b})\to a.b$. Si l'on identifie $\mathrm H^{2}(\mathrm A,\mathbf Z/2\,\mathbf Z)$ au sous-groupe du groupe de Brauer de $\mathrm A$ formé des $x$ tels que $2x=0$, l'élément $a.b$
\scanpage{37}
% archive sheet number: 36
\mspage{( 2 )}
% L1–2
s'interprète comme la classe de l'algèbre de quaternions définie par le couple $(a,b)$. Appliquant un théorème dû à Witt $(^{4})$, on a :

% L3–4
\textsc{Corollaire}. --- \emph{Le noyau de $\varphi$ est l'idéal engendré par les éléments $e_{a}+e_{b}-e_{a'}-e_{b'}$ tels que $e_{a+b}=e_{a'+b'}$ et $a.b=a'.b'$.}

% L5–11
Soit $\mathrm H$ un anneau gradué à degrés positifs, et $\widetilde{\mathrm H}=\mathbf Z\times(1+\widehat{\mathrm H}^{+})$ l'anneau filtré déduit de $\mathrm H$ par le procédé exposé par Grothendieck dans $(^{2})$. Alors la donnée d'un homomorphisme
\[
w^{1}:\quad\mathrm D(\mathrm A)\to\mathrm H^{1}
\]
tel que $w^{1}(e_{a})\,w^{1}(e_{b})=w^{1}(e_{a'})\,w^{1}(e_{b'})$ dès que $(a,b,a',b')$, vérifient les conditions du corollaire, définit un homomorphisme d'anneau filtré de $\mathrm Q(\mathrm A)$ dans $\widetilde{\mathrm H}$.

% L12–15
Prenons pour $\mathrm H$ l'algèbre de cohomologie $\mathrm H(\mathrm A)$ du groupe de Galois de $\mathrm A$ à valeurs dans $\mathbf Z/2\,\mathbf Z$, la multiplication étant définie par le cup-produit. Comme $\mathrm H^{1}(\mathrm A,\mathbf Z/2\,\mathbf Z)$ s'identifie à $\mathrm D(\mathrm A)$ on a une application
\[
w:\quad\mathrm Q(\mathrm A)\to\widetilde{\mathrm H(\mathrm A)}.
\]

% L16–17
L'homomorphisme $w$ définit donc des applications
\[
w^{i}:\quad\mathrm Q(\mathrm A)\to\mathrm H^{i}(\mathrm A,\mathbf Z/2\,\mathbf Z).
\]

% L18–21
Soient $(\mathrm E,\mathrm Q)$ un module quadratique et $q$ sa classe dans $\mathrm Q(\mathrm A)$; les $w^{i}(q)$ sont appelées les \emph{classes de Stiefel-Whitney} de $(\mathrm E,\mathrm Q)$ (ou de $q$). Comme $w$ est un homomorphisme d'anneaux, on a
\[
w^{n}(q+q')=\sum_{i+j=n}w^{i}(q)\,w^{j}(q').
\]

% L22–23
En particulier si $q$ est la classe du module $\left[\mathrm A^{p},\mathrm Q;\ \mathrm Q(x)=\displaystyle\sum_{1}^{p}a_{j}\,x_{j}^{2}\right]$, on a
\[
w^{1}(q)=\sum a_{j};\qquad w^{2}(q)=\sum_{i<j}a_{i}.a_{j}.
\]

% L24–31
On interprète $w^{1}(q)$ comme le discriminant de $(\mathrm E,\mathrm Q)$, et $w^{2}(q)$ de la façon suivante : si $rg\,(\mathrm E)=p$ est pair, soit $\mathrm C(q)$ (resp. $\mathrm C_{p}$) l'algèbre de Clifford de $(\mathrm E,\mathrm Q)$ [resp. de $p\,(\mathrm A,\mathrm I)$]. Il existe une extension galoisienne $\mathrm B$ de $\mathrm A$ telle que $(\mathrm E,\mathrm Q)$ et $p\,(\mathrm A,\mathrm I)$ soient isomorphes sur $\mathrm B$. Donc $\mathrm C(q)$ et $\mathrm C_{p}$ sont isomorphes sur $\mathrm B$. On en déduit un élément $u(q)$ du groupe de Brauer de $\mathrm A$ par le procédé habituel $(^{3})$, et l'on a $u(q)=w^{2}(q)$. Si $p$ est impair, on a la même interprétation en remplaçant $\mathrm C(q)$ par $\mathrm C^{+}(q)$ (resp. $\mathrm C_{p}$ par $\mathrm C_{p}^{+}$) $(^{1})$.

% L32–34
Enfin on vérifie que $w$ est un homomorphisme du foncteur $\mathrm Q(\mathrm A)$ dans le foncteur $\mathrm H(\mathrm A)$. Ceci provient du fait que $\mathrm H^{1}(\mathrm A,\mathbf Z/2\,\mathbf Z)$ est un foncteur covariant en $\mathrm A$.
\scanpage{38}
% archive sheet number: 37
\mspage{( 3 )}
% L1–5
\emph{Application aux corps de nombres algébriques.} --- Soit $\mathrm A$ un corps de nombres algébriques; à toute valeur absolue $v$ de $\mathrm A$, on associe le complété $\mathrm A_{v}$ de $\mathrm A$ pour cette valeur absolue. D'où une application
\[
\mathrm Q(\mathrm A)\to\prod\mathrm Q(\mathrm A_{v})
\]
et cette application est injective (théorème de Hasse).

% L6–9
Si $v$ est une valuation discrète, toute classe de forme sur $\mathrm A_{v}$ est déterminée par $\varepsilon$, $w^{1}$ et $w^{2}$ $(^{4})$. Si $\mathrm A_{v}=\mathbf R$, le corps des réels, alors $\mathrm H(\mathbf R)$ s'identifie à $\mathbf Z/2\,\mathbf Z\,[t]$. Toute classe de forme est définie par deux nombres entiers $m$ et $n$ ($m$, nombre de carrés positifs; $n$, nombre de carrés négatifs).

% L10–11
Posant $n=\displaystyle\sum n_{j}\,2^{j}$ $(n_{j}=0\text{ ou }1)$, on trouve
\[
w^{j}(q)=\binom{n}{j}\bmod 2;\qquad
\text{en particulier}\quad w^{2^{j}}=n_{j}.
\]

% L12–15
Donc une classe de forme sur $\mathbf R$ est complètement déterminée par $w$. De tout ce qui précède on déduit le diagramme commutatif suivant :
\diagram{\begin{tikzcd}[ampersand replacement=\&]
\mathrm Q(\mathrm A)\arrow[r]\arrow[d,"w"] \& \displaystyle\prod\mathrm Q(\mathrm A_{v})\arrow[d,"{(w_{v})}"]\\
\mathrm H(\mathrm A)\arrow[r] \& \displaystyle\prod\mathrm H(\mathrm A_{v})
\end{tikzcd}}

% L16–18
On démontre ainsi que l'application $w$ est injective : Sur un corps de nombres toute classe de forme quadratique est déterminée par son rang et ses classes de Stiefel-Whitney.

% L19–22
\emph{Remarque.} --- On ne sait pas si l'application
\[
\operatorname{grad}^{2}(\mathrm Q(\mathrm A))\to\mathrm H^{2}(\mathrm A,\mathbf Z/2\,\mathbf Z)
\]
est toujours injective, et l'on ne sait pas non plus si tout élément d'ordre $2$ du groupe de Brauer de $\mathrm A$ est somme d'éléments de la forme $a.b$.

% L23–26
{\small
$(^{1})$ N. \textsc{Bourbaki}, \emph{Algèbre}, chap. IX, Hermann, Paris, 1959.\\
$(^{2})$ A. \textsc{Grothendieck}, \emph{Bull. Soc. math. Fr.}, 86, 1958.\\
$(^{3})$ T. A. \textsc{Springer}, \emph{Akad. Amst. Proc.}, 62 A, 1959.\\
$(^{4})$ E. \textsc{Witt}, \emph{Journal de Crelle}, 176, 1937.\par}
\scanpage{39}
% archive sheet number: 38
% L1–2
\begin{center}
Extrait des \emph{Comptes rendus des séances de l'Académie des Sciences},\\
t. 255, p. 1366-1368, séance du 10 septembre 1962.

% L3–7
GAUTHIER-VILLARS \& C\textsuperscript{ie},\\
55, Quai des Grands-Augustins, Paris (6\textsuperscript{e}),\\
Éditeur-Imprimeur-Libraire.\\
162396

Imprimé en France.
\end{center}
